\subsection*{升级专题：已经成功的 O1 怎样补建通知}

前两节已经让新系统完成接收身份 I 与效果身份 B 的共同仲裁。现在把源端与消费端一起放回升级现场：第八章已有成功 O1，旧 R 没有 Outbox；旧服务还可能发送过另一种格式的通知。上线第九章的新 R，只会改变后来执行的事务，已经结束的成功也要有明确的迁移位置。

为这批旧成功记录安排一次受控补建。切换时先停止旧 R 入口，等待在途本地事务结束，再把当时已经成功的 C 及其原事实摘要登记为持久的旧记录集合。这个集合也称为 cohort，即本次切换明确圈定的一批记录。登记由迁移程序保存，补建者据它判断哪些成功事实属于旧规则。全体补建完成并核对以后，才开放只走新规则的 R。原结果的只读查询可以继续，尚未决的 C 后来取得成功时直接走新 R，平台付款身份保持原值。暂停、登记与补建屏障共同划出可核查的边界，每项旧成功都有明确的补建归属。

\CTenFigure{ch09-migration-barrier}{先圈定旧成功集合并完成补建，再开放新 R，形成可恢复的源端切换屏障}{fig:ch09-migration-barrier}

补建事务按 C1→T1 锁定原关联，读取旧集合中的登记，确认 C1 已成功、P1 已 \code{paid}、T1 已 \code{done}，并核对唯一 O1 与保存结果中的 P/C/O/X、币种和金额一致。然后从原购买输入、报价快照和已保存支付事实构造通知，按同一事实唯一约束查找事件。尚无事件才插入，已有事件则逐字段比较并恢复。事件与本项补建完成共同提交。并发补建者在原 C 处串行，退出以后按登记继续。零行、冲突或原事实摘要不符，均整笔回滚并登记核查。补建只增加发布依据，原付款、订单、结果和任务终态保持原值。

若旧服务曾直接发送过通知，补建后的消息可能再次报告 O1 已付款。新分配的事件编号只能统一新路径的副本。旧副本若没有这个编号，下游就要沿 O1 对应的履约动作认出相同效果。启用这批补发以前，旧、新通知应进入同一项履约效果约束。对于发送情况不明的旧通知，迁移记录保存未知结论和负责核查的归属。

新规则启用后，受控补建只面向已登记的旧成功集合。已经按新 R 提交的成功记录若缺少事件，普通结果重放报告完整性异常，并保留核查归属。核查者据原提交及发送证据判断是否误删、是否已有消息到达，再决定恢复办法。普通重放径直分配一个新事件，会把这些情况混在一起。
